\documentclass[10pt,a4paper]{amsart}
\usepackage{amsmath,amssymb,amsthm,graphicx}
% Retained minimal preamble for the Stacks Project chapter more-morphisms.tex (More on Morphisms).
% Verbatim excerpt of the chapter preamble: Xy-pic setup (native preamble lines 14-18) and the native
% theorem declarations with the mathematical macros (native preamble lines 154-217). Copied mechanically;
% nothing added, removed or redefined. The wrapper class and the amsmath/amssymb/amsthm/graphicx base
% packages are supplied by the assembler, so the class/externaldocument part of the native preamble is omitted.
% For commutative diagrams we use Xy-pic
\usepackage[all]{xy}

% We use 2cell for 2-commutative diagrams.
\xyoption{2cell}
% Theorem environments.
%
\theoremstyle{plain}
\newtheorem{theorem}[subsection]{Theorem}
\newtheorem{proposition}[subsection]{Proposition}
\newtheorem{lemma}[subsection]{Lemma}

\theoremstyle{definition}
\newtheorem{definition}[subsection]{Definition}
\newtheorem{example}[subsection]{Example}
\newtheorem{exercise}[subsection]{Exercise}
\newtheorem{situation}[subsection]{Situation}

\theoremstyle{remark}
\newtheorem{remark}[subsection]{Remark}
\newtheorem{remarks}[subsection]{Remarks}

\numberwithin{equation}{subsection}

% Macros
%
\def\lim{\mathop{\mathrm{lim}}\nolimits}
\def\colim{\mathop{\mathrm{colim}}\nolimits}
\def\Spec{\mathop{\mathrm{Spec}}}
\def\Hom{\mathop{\mathrm{Hom}}\nolimits}
\def\Ext{\mathop{\mathrm{Ext}}\nolimits}
\def\SheafHom{\mathop{\mathcal{H}\!\mathit{om}}\nolimits}
\def\SheafExt{\mathop{\mathcal{E}\!\mathit{xt}}\nolimits}
\def\Sch{\mathit{Sch}}
\def\Mor{\mathop{\mathrm{Mor}}\nolimits}
\def\Ob{\mathop{\mathrm{Ob}}\nolimits}
\def\Sh{\mathop{\mathit{Sh}}\nolimits}
\def\NL{\mathop{N\!L}\nolimits}
\def\CH{\mathop{\mathrm{CH}}\nolimits}
\def\proetale{{pro\text{-}\acute{e}tale}}
\def\etale{{\acute{e}tale}}
\def\QCoh{\mathit{QCoh}}
\def\Ker{\mathop{\mathrm{Ker}}}
\def\Im{\mathop{\mathrm{Im}}}
\def\Coker{\mathop{\mathrm{Coker}}}
\def\Coim{\mathop{\mathrm{Coim}}}

% Boxtimes
%
\DeclareMathSymbol{\boxtimes}{\mathbin}{AMSa}{"02}

%
% Macros for moduli stacks/spaces
%
\def\QCohstack{\mathcal{QC}\!\mathit{oh}}
\def\Cohstack{\mathcal{C}\!\mathit{oh}}
\def\Spacesstack{\mathcal{S}\!\mathit{paces}}
\def\Quotfunctor{\mathrm{Quot}}
\def\Hilbfunctor{\mathrm{Hilb}}
\def\Curvesstack{\mathcal{C}\!\mathit{urves}}
\def\Polarizedstack{\mathcal{P}\!\mathit{olarized}}
\def\Complexesstack{\mathcal{C}\!\mathit{omplexes}}
% \Pic is the operator that assigns to X its picard group, usage \Pic(X)
% \Picardstack_{X/B} denotes the Picard stack of X over B
% \Picardfunctor_{X/B} denotes the Picard functor of X over B
\def\Pic{\mathop{\mathrm{Pic}}\nolimits}
\def\Picardstack{\mathcal{P}\!\mathit{ic}}
\def\Picardfunctor{\mathrm{Pic}}
\def\Deformationcategory{\mathcal{D}\!\mathit{ef}}
\usepackage{cleveref}
\crefname{originallemma}{original lemma}{original lemmas}
\Crefname{originallemma}{Original lemma}{Original lemmas}
\makeatletter
\newlabel{algebra-lemma-helper-polynomial}{{algebra-lemma-helper-polynomial}{source}{}{sourcepointer-1}{}}
\newlabel{algebra-lemma-helper-polynomial@cref}{{[originallemma][0][]algebra-lemma-helper-polynomial}{[1][0][]source}{}{}{}}
\makeatother
\crefname{originallemma}{original lemma}{original lemmas}
\Crefname{originallemma}{Original lemma}{Original lemmas}
\makeatletter
\newlabel{algebra-theorem-nullstellensatz}{{algebra-theorem-nullstellensatz}{source}{}{sourcepointer-2}{}}
\newlabel{algebra-theorem-nullstellensatz@cref}{{[originallemma][0][]algebra-theorem-nullstellensatz}{[1][0][]source}{}{}{}}
\makeatother
\title{An irreducible polynomial over a domain has geometrically irreducible fibres over a dense open}
\author{The Stacks Project Authors}
\date{stacks-project:more-morphisms.tex:local-snapshot-20261011; GNU FDL 1.2 or later}
\begin{document}
\maketitle
\noindent\textit{Editorial note.} Statement span is lines 7491-7500 of sources/stacks-more-morphisms/pinned-original-source.tex; proof span is lines 7502-7554, the complete author proof (51 body lines).\par\medskip
\noindent\textit{Editorial note.} Navigation stubs: 2 cited-but-not-included labels are supplied as non-load-bearing pointers with their pinned tags where the snapshot tag index has them.\par\medskip
\noindent\textit{Editorial note.} extra\_packages is empty on purpose: the retained preamble already loads xy with options and a second declaration clashes in the assembler wrapper.\par\medskip
\noindent\textit{Editorial note.} Text is retained verbatim, including the author's own typos and placeholders.\par\medskip
\noindent\textit{Allowed standard prerequisites (editorial).}\begin{itemize}
\item Explicit prerequisite boundary. Irreducibility of polynomials over a field, localisation, and the dictionary between fibres and base change; the reduction is retained.
\end{itemize}\medskip
\ifdefined\hypertarget\hypertarget{sourcepointer-1}{}\fi
\noindent\textit{Source reference.} Original lemma algebra-lemma-helper-polynomial: sources/stacks-more-morphisms/pinned-original-source.tex, source lines 7554--7554. Cited by the retained statement or proof but not part of the collected spans; supplied only as a navigation pointer. The argument does not depend on its content.\par\medskip
\ifdefined\hypertarget\hypertarget{sourcepointer-2}{}\fi
\noindent\textit{Source reference.} Original lemma algebra-theorem-nullstellensatz: sources/stacks-more-morphisms/pinned-original-source.tex, source lines 7554--7554. Cited by the retained statement or proof but not part of the collected spans; supplied only as a navigation pointer. The argument does not depend on its content.\par\medskip
\begin{lemma}
\label{lemma-irreducible-polynomial-over-domain}
Let $A$ be a domain with fraction field $K$.
Let $P \in A[x_1, \ldots, x_n]$.
Denote $\overline{K}$ the algebraic closure of $K$.
Assume $P$ is irreducible in $\overline{K}[x_1, \ldots, x_n]$.
Then there exists a $f \in A$ such that
$P^\varphi \in \kappa[x_1, \ldots, x_n]$ is irreducible for all
homomorphisms $\varphi : A_f \to \kappa$ into fields.
\end{lemma}

\begin{proof}
There exists an automorphism $\Psi$ of $A[x_1, \ldots, x_n]$ over $A$
such that $\Psi(P) = ax_n^d +$ lower order terms in $x_n$ with
$a \not = 0$, see
Algebra, Lemma \ref{algebra-lemma-helper-polynomial}.
We may replace $P$ by $\Psi(P)$ and we may replace $A$ by $A_a$.
Thus we may assume that $P$ is monic in $x_n$ of degree $d > 0$.
For $i = 1, \ldots, n - 1$ let $d_i$ be the degree of $P$ in $x_i$.
Note that this implies that $P^\varphi$ is monic of degree $d$ in $x_n$
and has degree $\leq d_i$ in $x_i$ for every homomorphism
$\varphi : A \to \kappa$ where $\kappa$ is a field.
Thus if $P^\varphi$ is reducible, then we can write
$$
P^\varphi = Q_1 Q_2
$$
with $Q_1, Q_2$ monic of degree $e_1, e_2 \geq 0$ in $x_n$ with
$e_1 + e_2 = d$ and having degree $\leq d_i$ in $x_i$ for
$i = 1, \ldots, n - 1$. In other words we can write
\begin{equation}
\label{equation-factors}
Q_j = x_n^{e_j} + \sum\nolimits_{0 \leq l < e_j}
\left( \sum\nolimits_{L \in \mathcal{L}} a_{j, l, L} x^L \right) x_n^l
\end{equation}
where the sum is over the set $\mathcal{L}$ of multi-indices $L$
of the form $L = (l_1, \ldots, l_{n - 1})$ with $0 \leq l_i \leq d_i$.
For any $e_1, e_2 \geq 0$ with $e_1 + e_2 = d$ we consider the $A$-algebra
$$
B_{e_1, e_2} =
A[\{a_{1, l, L}\}_{0 \leq l < e_1, L \in \mathcal{L}},
\{a_{2, l, L}\}_{0 \leq l < e_2, L \in \mathcal{L}}]/(\text{relations})
$$
where the $(\text{relations})$ is the ideal generated by the coefficients
of the polynomial
$$
P - Q_1Q_2 \in
A[\{a_{1, l, L}\}_{0 \leq l < e_1, L \in \mathcal{L}},
\{a_{2, l, L}\}_{0 \leq l < e_2, L \in \mathcal{L}}][x_1, \ldots, x_n]
$$
with $Q_1$ and $Q_2$ defined as in (\ref{equation-factors}). OK, and
the assumption that $P$ is irreducible over $\overline{K}$ implies that
there does not exist any $A$-algebra homomorphism
$B_{e_1, e_2} \to \overline{K}$. By the Hilbert Nullstellensatz, see
Algebra, Theorem \ref{algebra-theorem-nullstellensatz}
this means that $B_{e_1, e_2} \otimes_A K = 0$.
As $B_{e_1, e_2}$ is a finitely generated $A$-algebra this signifies that
we can find an $f_{e_1, e_2} \in A$ such that
$(B_{e_1, e_2})_{f_{e_1, e_2}} = 0$. By construction this means that
if $\varphi : A_{f_{e_1, e_2}} \to \kappa$ is a homomorphism to a field,
then $P^\varphi$ does not have a factorization $P^\varphi = Q_1 Q_2$
with $Q_1$ of degree $e_1$ in $x_n$ and $Q_2$ of degree $e_2$ in $x_n$.
Thus taking
$f = \prod_{e_1, e_2 \geq 0, e_1 + e_2 = d} f_{e_1, e_2}$ we win.
\end{proof}

\end{document}
